%-------------------------------------------------------------
%  CV Radu-Lucian Moldovan -- versiune cu cuvinte cheie ATS
%  Compilează în Overleaf cu pdfLaTeX (setarea implicită)
%
%  >>> ALEGE ȚINTA AICI (modifici doar cuvântul dintre acolade):
%      android   -> Junior Android / Mobile Developer
%      embedded  -> Junior Embedded / Automation Engineer
%      software  -> Junior Software Developer (general)
%-------------------------------------------------------------
\providecommand{\TINTA}{software}

\documentclass[10pt,a4paper]{article}

\usepackage[utf8]{inputenc}
\usepackage[T1]{fontenc}
\usepackage[scaled]{helvet}
\renewcommand{\familydefault}{\sfdefault}
\usepackage{combelow}
\usepackage{newunicodechar}
\newunicodechar{ș}{\cb{s}}
\newunicodechar{Ș}{\cb{S}}
\newunicodechar{ț}{\cb{t}}
\newunicodechar{Ț}{\cb{T}}

\usepackage[margin=1.4cm,top=1.0cm,bottom=1.0cm]{geometry}
\usepackage{xcolor}
\usepackage{titlesec}
\usepackage{enumitem}
\usepackage{tabularx}
\usepackage{etoolbox}
\usepackage[hidelinks]{hyperref}

\definecolor{accent}{HTML}{1F4E79}
\definecolor{muted}{HTML}{555555}

\pagestyle{empty}
\setlength{\parindent}{0pt}

\titleformat{\section}{\large\bfseries\color{accent}}{}{0pt}{\MakeUppercase}[\vspace{-6pt}{\color{accent}\rule{\linewidth}{0.6pt}}]
\titlespacing*{\section}{0pt}{7pt}{4pt}

\setlist[itemize]{leftmargin=1.2em,itemsep=1pt,topsep=2pt,parsep=0pt,label=\textbullet}

% \entry{Titlu}{Perioadă}{Organizație}{Locație}
\newcommand{\entry}[4]{%
  \textbf{#1} \hfill {\color{muted}#2}\\
  \textit{#3} \hfill {\color{muted}\small #4}\par\vspace{2pt}}

% \project{Nume}{Tehnologii}{Link}
\newcommand{\project}[3]{%
  \textbf{#1} \hfill {\small\href{https://#3}{#3}}\\
  {\small\color{muted}\textit{#2}}\par\vspace{1pt}}

% \doar{android,software}{text} -> apare doar pentru țintele din listă
\newtoggle{potrivit}
\newcommand{\verificatinta}[1]{\ifdefstring{\TINTA}{#1}{\toggletrue{potrivit}}{}}
\newcommand{\doar}[2]{\togglefalse{potrivit}\forcsvlist{\verificatinta}{#1}\iftoggle{potrivit}{#2}{}}

%================ BLOCURI DE CONȚINUT =================

\newcommand{\proiectHarmony}{%
\project{Harmony -- Player muzical Android}%
  {\doar{android}{Kotlin, Jetpack Compose, MVVM, Hilt, Room, Coroutines, Media3/ExoPlayer, GitHub Actions}%
   \doar{embedded}{Kotlin, C/C++ (Android NDK, CMake), DSP audio, TCP sockets, programare concurentă}%
   \doar{software}{Kotlin, SQL (Room/SQLite), REST API, TCP sockets, Git, CI/CD (GitHub Actions)}}%
  {github.com/R4duQ/HarmonyApp}
\begin{itemize}
  \item Aplicație construită de la zero, cu \textbf{arhitectură multi-modul (MVVM)}, injecție de dependențe (Hilt)
        și bază de date locală \textbf{SQL (Room)} cu căutare full-text (FTS4) incrementală.
  \doar{android}{%
  \item UI în \textbf{Jetpack Compose} cu animații, gesturi și suport landscape, proiectat întâi în Figma;
        playback în fundal cu \textbf{Media3/ExoPlayer} și persistența stării.
  \item Integrare \textbf{Android Auto} (MediaSession), recuperare automată din erori de playback (retry, circuit breaker).}
  \doar{embedded}{%
  \item \textbf{Procesare de semnal (DSP) în timp real}: egalizator, ReplayGain și limiter soft pe PCM 16-bit,
        cu un modul nativ \textbf{C/C++} compilat prin Android NDK și CMake.
  \item Rezolvarea unor \textbf{race conditions} în procesorul audio și crossfade cu două motoare de redare sincronizate.}
  \doar{android,embedded}{%
  \item \textbf{Motor audio custom}: playback lossless (FLAC), egalizator, ReplayGain și crossfade.}
  \doar{software}{%
  \item \textbf{Integrare REST API} (serviciu public de metadate) pentru preluarea coperților la runtime.
  \item \textbf{Debugging} de concurență (race conditions) și recuperare automată din erori (retry, circuit breaker).}
  \item Funcție de descărcare prin \textbf{Soulseek}: client \textbf{peer-to-peer} propriu peste TCP sockets, cu
        gestionarea conexiunilor concurente, cozi de descărcare, timeout-uri, cache LRU și deduplicarea căutărilor.
  \item Algoritm \textbf{Smart Shuffle} bazat pe scor de similaritate acustică între piese.
  \item Release-uri semnate automat prin \textbf{CI/CD (GitHub Actions)}.
\end{itemize}}

\newcommand{\proiectSemafor}{%
\textbf{Semafor inteligent} \textbar\ {\small\textit{Arduino, C/C++, senzori}}\par\vspace{1pt}
\begin{itemize}
  \item Sistem de semaforizare care își adaptează ciclul pe baza datelor de la \textbf{senzori};
        proiect de electronică, automatizare și \textbf{programare embedded} pe microcontroler.
\end{itemize}}

% Rânduri pentru tabelul de competențe
\newcommand{\rLimbajeMobile}{\textbf{Limbaje:} & Kotlin, Java, C, C++, C\#, Python, Go, Lua, SQL, Objective-C \\}
\newcommand{\rLimbajeEmbedded}{\textbf{Limbaje:} & C, C++, Python, Lua, Kotlin, Java, C\#, Go, SQL \\}
\newcommand{\rLimbajeSoftware}{\textbf{Limbaje:} & Java, Kotlin, C\#, Python, C, C++, Go, Lua, SQL \\}
\newcommand{\rAndroid}{\textbf{Android:} & Jetpack Compose, MVVM, Hilt (Dependency Injection), Room, Coroutines, Media3/ExoPlayer, Android Auto, Android NDK \\}
\newcommand{\rEmbedded}{\textbf{Embedded:} & Arduino, microcontrolere, senzori și actuatoare, DSP, reglare automată (PID), IoT \\}
\newcommand{\rConcepte}{\textbf{Concepte:} & OOP, structuri de date și algoritmi, programare concurentă, REST API, baze de date relaționale, debugging \\}
\newcommand{\rUnelte}{\textbf{Unelte:} & Git, GitHub, GitHub Actions (CI/CD), Gradle, CMake, Android Studio, Figma, Google Colab \\}
\newcommand{\rLimbi}{\textbf{Limbi străine:} & Română (nativ), Engleză (C1) \\}
\newcommand{\rPermis}{\textbf{Permis:} & Categoriile AM și B \\}

%=====================================================

\begin{document}

%---------------- ANTET ----------------
\begin{center}
  {\Huge\bfseries Radu-Lucian Moldovan}\\[4pt]
  {\large\color{accent} Student Ingineria Sistemelor}\\[5pt]
  \small
  Timișoara, România \;\textbar\;
  \href{mailto:radu.moldovan2005@gmail.com}{radu.moldovan2005@gmail.com} \;\textbar\;
  +40 732 815 096 \;\textbar\;
  \href{https://github.com/R4duQ}{github.com/R4duQ}
\end{center}

%---------------- PROFIL ----------------
\section{Profil}
\doar{android}{%
Student în anul III la Ingineria Sistemelor (UPT), orientat spre \textbf{dezvoltare Android nativă} în
\textbf{Kotlin} și \textbf{Jetpack Compose}. Am proiectat și dezvoltat de la zero \textbf{Harmony}, un player
muzical complet, publicat pe GitHub, cu arhitectură \textbf{MVVM multi-modul}, Hilt, Room, Coroutines,
Media3/ExoPlayer, motor audio propriu și integrare Android Auto. Livrez release-uri automate prin
\textbf{CI/CD (GitHub Actions)} și urmăresc performanța, stabilitatea și experiența utilizatorului, de la
design în Figma până la debugging de concurență.}
\doar{embedded}{%
Student în anul III la Ingineria Sistemelor (UPT), cu pregătire în \textbf{sisteme embedded},
reglare automată și \textbf{procesarea semnalelor}. Am experiență practică în \textbf{C/C++} pe
microcontrolere (Arduino, senzori și actuatoare) și în \textbf{DSP audio în timp real}: egalizator,
ReplayGain și limiter pe PCM 16-bit, cu cod nativ integrat prin Android NDK și CMake. Depanez metodic
probleme de concurență și sincronizare și duc proiectele de la prototip până la release.
Caut o poziție de \textbf{Junior Embedded Software Engineer}.}
\doar{software}{%
Student în anul III la Ingineria Sistemelor (UPT), cu bază solidă în \textbf{programare, OOP, structuri
de date și algoritmi} și experiență în mai multe limbaje (Kotlin, Java, C/C++, C\#, Python, Go).
Am proiectat și livrat de la zero o aplicație complexă, publicată pe GitHub, cu arhitectură modulară,
bază de date \textbf{SQL}, integrare \textbf{REST API}, client de rețea peer-to-peer și pipeline \textbf{CI/CD}.
Învăț rapid, depanez metodic probleme dificile (concurență, erori de rețea) și duc proiectele până la
release. Caut o poziție de \textbf{Junior Software Developer}.}

%---------------- EDUCAȚIE ----------------
\section{Educație}
\entry{Ingineria Sistemelor -- Licență}{2024 -- prezent}
      {Universitatea Politehnica Timișoara, Facultatea de Automatică și Calculatoare}{Timișoara}
\begin{itemize}
  \item Programare și structuri de date: C, C++, Java, Python, algoritmi
  \item Sisteme embedded și microcontrolere: achiziție de date, senzori, actuatoare
  \item Modelare matematică, identificarea sistemelor, reglare automată (PID)
  \item Rețele industriale, protocoale în timp real / IoT, prelucrarea semnalelor și a imaginilor
\end{itemize}
\vspace{4pt}

\entry{Matematică-Informatică, intensiv Informatică}{2020 -- 2024}
      {Colegiul Național de Informatică ,,Carmen Sylva''}{Petroșani}

%---------------- PROIECTE ----------------
\section{Proiecte}
\ifdefstring{\TINTA}{embedded}
  {\proiectSemafor\vspace{4pt}\proiectHarmony}
  {\proiectHarmony\vspace{4pt}\proiectSemafor}

%---------------- VOLUNTARIAT ----------------
\section{Voluntariat}
\entry{Voluntar, Departamentul Educațional}{2024 -- prezent}{Liga AC}{Timișoara}
\begin{itemize}
  \item Implicat în planificarea și coordonarea evenimentelor \textbf{UniHack} și \textbf{ITEC}.
\end{itemize}
\vspace{4pt}

\entry{Participant, proiect Erasmus+}{2020 -- 2022}{Colegiul Național de Informatică ,,Carmen Sylva'', Petroșani}{România / Turcia}
\begin{itemize}
  \item Schimb de experiență în \textbf{Fethiye, Turcia}, alături de elevi din alte țări partenere.
  \item Implicat în organizarea activităților proiectului la liceu, în România.
\end{itemize}

%---------------- COMPETENȚE ----------------
\section{Competențe}
\ifdefstring{\TINTA}{embedded}{%
\begin{tabularx}{\linewidth}{@{}l X@{}}
  \rLimbajeEmbedded \rEmbedded \rConcepte \rUnelte \rAndroid \rLimbi \rPermis
\end{tabularx}}{}%
\ifdefstring{\TINTA}{software}{%
\begin{tabularx}{\linewidth}{@{}l X@{}}
  \rLimbajeSoftware \rConcepte \rUnelte \rAndroid \rEmbedded \rLimbi \rPermis
\end{tabularx}}{}%
\ifdefstring{\TINTA}{android}{%
\begin{tabularx}{\linewidth}{@{}l X@{}}
  \rLimbajeMobile \rAndroid \rConcepte \rUnelte \rEmbedded \rLimbi \rPermis
\end{tabularx}}{}

\end{document}
